\documentclass[10pt,a4paper]{amsart}
\usepackage[margin=19mm]{geometry}
\usepackage{amsmath,amssymb,mathtools}
\usepackage[scr=rsfs,cal=euler]{mathalfa}
\usepackage{xfrac}
\usepackage[unicode,hidelinks,bookmarks=false]{hyperref}
\newcommand{\ie}{i.e.\ }
\newcommand{\cf}{cf.\ }
\newcommand{\resp}{respectively\ }
\newcommand{\Subsection}{\subsection}
\providecommand{\nobreakdash}{\nobreak\hbox{-}}
\setlength{\emergencystretch}{2em}
\multlinegap0pt
\usepackage{bm}
\usepackage{bbm}
\usepackage{dsfont}
\usepackage{xspace}
\usepackage{cancel}
\usepackage{mathtools}
\usepackage{amsthm}
\makeatletter
\@ifundefined{theorem}{\newtheorem{theorem}{Theorem}}{}
\@ifundefined{lemma}{\newtheorem{lemma}[theorem]{Lemma}}{}
\makeatother
\title[Clarification and Coinduction of Tambara Functors]{Clarification and Coinduction of Tambara Functors}
\author{Noah Wisdom}
\date{Source: arXiv:2505.08066v2 (CC BY 4.0)}
\begin{document}
\maketitle

\noindent\emph{Source and licence.} Clarification and Coinduction of Tambara Functors. Version arXiv:2505.08066v2 (CC BY 4.0), distributed under the Creative Commons Attribution 4.0 International licence, as recorded in the arXiv licence metadata for this exact version and shown on the abstract page https://arxiv.org/abs/2505.08066v2. Full rights notice: licenses/arxiv-2505.08066-CC-BY-4.0.txt.\par\smallskip

\noindent\textit{Editorial scope.} Counted unit: the Lemma whose source label is \texttt{lem:norm-of-idempotents-is-additive} in the article (source0.tex lines 461--463, source label \texttt{lem:norm-of-idempotents-is-additive}), delivered together with the article's own complete original proof of it (source0.tex lines 465--484, one \texttt{proof} environment of 20 lines). No mathematics is written, repaired or re-derived: statement and proof are the article's own text line for line. Required context retained uncounted: none. Every definition, display and label of those lines is retained unchanged. Omitted from this package and not redistributed: 1--460, 464--464, 485--1151. Nothing inside a retained excerpt is omitted or altered: every retained body is delivered exactly as the article's lines are written, so no omission receipt of any kind accompanies this package. Citations: the retained excerpts carry no citation command at all, so no bibliography entry of the article is transcribed and no reference is redistributed. Reference adaptation: none was needed and none was made; every reference inside the delivered unit resolves inside it, so no unresolved reference or citation is produced. Printed slips kept literal: no printed word, symbol, hypothesis or number of the counted statement or of its proof was altered. Layout and class: the article's own class is \texttt{amsart} with packages including xy, graphicx, amsmath, amssymb, amsthm, amsfonts, eucal, hyperref, comment, cleveref, microtype, tikz-cd, spectralsequences, todonotes; the wrapper is the lane's pinned \texttt{amsart} editorial wrapper at 10pt on A4, which restates the theorem-like environments the retained text needs and adds the article's own macro definitions that the retained text needs (none), with extra wrapper packages (bm, bbm, dsfont, xspace, cancel, mathtools). The delivered numbering is the unit's own. Assets: no retained span contains a figure environment, a graphics-inclusion command, a PDF-page inclusion, a table or a TikZ picture; no image file is copied into this package, the package declares no asset dependency and no image file is redistributed with it. Licence: version arXiv:2505.08066v2 of this article is distributed under the Creative Commons Attribution 4.0 International licence, as recorded in the arXiv licence metadata for this exact version and shown on the abstract page https://arxiv.org/abs/2505.08066v2. A full rights notice is delivered at licenses/arxiv-2505.08066-CC-BY-4.0.txt. Characters: every retained excerpt is pure ASCII, so the delivered engine, XeLaTeX with its default Latin Modern text font, reports no missing character.
\begin{lemma}\label{lem:norm-of-idempotents-is-additive}
	Let $a$ and $b$ be mutually orthogonal idempotents in $R(G/e)$ which are fixed by $G$, and let $H$ be any subgroup of $G$. Then $\mathrm{Nm}_e^H(a+b) = \mathrm{Nm}_e^H(a)+\mathrm{Nm}_e^H(b)$.
\end{lemma}

\begin{proof}
	The exponential formula gives an alternate description of the norm of a sum. We have 
    \[ 
        \Pi_f (G/e \sqcup G/e) \cong G/H \times \mathrm{Fun}(G/H,\{ \mathrm{l} , \mathrm{r} \}) 
    \] 
    as $G$-sets.
	
	Observe the isomorphism $G/e \times_{G/H} \Pi_f (G/e \sqcup G/e) \cong G/e \times \mathrm{Fun}(G/H,\{ \mathrm{l} , \mathrm{r} \})$, which is a disjoint union of copies of $G/e$, as $G$ acts diagonally. Each orbit has a unique representative given by the tuple with $e$ in the first coordinate. The restriction is a product of projections: it is a map \[ R(G/e) \times R(G/e) \rightarrow \prod_{(e,\sigma : H \rightarrow \{ \mathrm{l} , \mathrm{r} \})} R(G/e) \] where the map $R(G/e) \times R(G/e) \rightarrow R(G/e)$ determining the $(e,\sigma)$th factor is projection onto the left factor in the domain if $\sigma(e) = \mathrm{l}$, or the right factor if $\sigma(e) = \mathrm{r}$, followed by possibly postcomposing with left multiplication by an element of $G$. Since $a$ and $b$ are fixed by $G$ by assumption, the restriction sends $(a,b)$ to the element in the product $\Pi_{(e,\sigma : H \rightarrow \{ \mathrm{l} , \mathrm{r} \})} R(G/e)$ whose value in the $(e,\sigma)$th coordinate is $a$ if $\sigma(e) = \mathrm{l}$ and $b$ if $\sigma(e) = \mathrm{r}$.
	
	The norm along the projection onto $\Pi_f (G/e \sqcup G/e)$ is the product of a collection of norm and multiplication maps. Two factors in this product description respectively send the elements $a$ and $b$ in the factors of $R(G/e)$ corresponding to the two constant functions $\sigma$ to $\mathrm{Nm}_e^H(a)$ and $\mathrm{Nm}_e^H(b)$. To conclude the proof, it suffices to show that every other coordinate of the image of the norm along $G/e \times_{G/H} \Pi_f ( G/e \sqcup G/e) \rightarrow \Pi_f (G/e \sqcup G/e)$ is zero.
	
	The quotient map along which we are taking the norm has the form \[ (g, \sigma : gH \rightarrow \{ \mathrm{l} , \mathrm{r} \}) \mapsto (gH, \sigma : gH \rightarrow \{ \mathrm{l} , \mathrm{r} \}) . \] If $\sigma : gH \rightarrow \{ \mathrm{l} , \mathrm{r} \}$ is not constant, it suffices to find two elements $(e, \sigma_i : H \rightarrow \{ \mathrm{l} , \mathrm{r} \})$ in the preimage of the $G$-orbit of $(H, \sigma : H \rightarrow \{ \mathrm{l} , \mathrm{r} \})$ which have the property $\sigma_1(e) \neq \sigma_2(e)$, because this corresponds to the value of the norm at the factor $R(G \cdot (H,\sigma))$ of $R(\Pi_f (G/e \sqcup G/e))$ being divisible by \[ \mathrm{Nm}_e^K(a) \mathrm{Nm}_e^K(b) = \mathrm{Nm}_e^K(ab) = \mathrm{Nm}_e^K(0) = 0 \] where $K$ is the stabilizer of $(H,\sigma)$. In other words, the composition 
    \begin{align*}
        R(G/e) \times R(G/e) & \rightarrow R(G/e \times_{G/H} \Pi_f (G/e \sqcup G/e)) \\
        & \rightarrow R(\Pi_f (G/e \sqcup G/e)) \rightarrow R(G \cdot (H,\sigma))
    \end{align*} 
    sends $(a,b)$ to zero.
	
	Since $\sigma$ is not constant, there exists $h \in H$ such that $\sigma(h) \neq \sigma(1)$. Now $(e, \sigma)$ and $(h^{-1},\sigma)$ are two distinct preimages of $(H,\sigma)$, and $h(h^{-1},\sigma) = (e,x \mapsto \sigma(hx))$ implies $(e,x\mapsto \sigma(hx))$ is in the preimage of the $G$-orbit of $(H, \sigma : H \rightarrow \{ \mathrm{l} , \mathrm{r} \})$. Therefore $(e,\sigma)$ and $(e,x \mapsto \sigma(hx))$ have the desired property.
\end{proof}

\end{document}
